% Lösungen Zone (zusammenbau v0.1)
\einheitenkopf[][Kennst du schon]{Das kennst du schon}

Zuordnung: 1 → die Potenzregel in Einheit 1 hängt am Unterschied zwischen Vorfaktor und Exponent; negative und rationale Exponenten für den Vorrat · 2 → jeder Nachweis in vorgegebener Form in Einheit 1 und 3 · 3 → die fhr-Terme tragen fast alle solche Koeffizienten, Einheit 1 · 4 → Grundlage von Einheit 2 und 3 und des Ausklammerns · 5 → für die Graphenaufgaben in Einheit 2 und 3 (Werte am Graphen ablesen, die Steigung null am Tiefpunkt erkennen) · 6 → die Anschlussleistungen der fhr-Zeilen (Punkt auf dem Graphen, Anstieg im Achsenschnittpunkt, Anstiege vergleichen) · 7 → nur für die Prüfungshöhe von Einheit 2 (hundertste Ableitung als Verschiebung) · 8 → die Potenzregel in Einheit 1 hängt am Unterschied zwischen Vorfaktor und Exponent; negative und rationale Exponenten für den Vorrat · 9 → die Potenzregel in Einheit 1 hängt am Unterschied zwischen Vorfaktor und Exponent; negative und rationale Exponenten für den Vorrat

\erg{1}{a) $x^8$ \quad b) $x^{-5}$}
\erg{2}{a) $4x + 20$ \quad b) $(5x - 2x^2) \cdot e^x$}
\erg{3}{a) $2$ \quad b) $\frac{3}{2} = 1{,}5$}
\erg{4}{a) $1$ \quad b) $e^{-2} \approx 0{,}14$}
\erg{5}{a) Steigung = $1$ \quad b) $-1{,}5$: der Graph fällt dort, pro Schritt nach rechts um $1{,}5$ nach unten.}
\erg{6}{a) $f(3) = 13$ \quad b) $f'(-2) = 8 - 5 = 3$}
\erg{7}{a) $4$ nach rechts verschoben \quad b) $2{,}5$ nach links verschoben}
\erg{8}{a) Der Faktor $2$ wurde nicht mit hoch $4$ genommen; richtig: $(2x)^4 = 16x^4$.}
\erg{9}{a) $(5x)^3 = 125x^3$}
